\section{Verification and Validation}

\subsection{Unit Tests}

The generator project contains an extensive suite of unit tests located in
\texttt{tests/generator/tests/}. In total \emph{73} tests exercise individual
modules such as the model classes, rule engine, parsing utilities and test
strategies.  Model tests validate the behaviour of \texttt{Variable},
\texttt{State}, \texttt{Action} and \texttt{Machine} objects. They ensure that
states correctly track actions, variable mappings are applied, and rule checking
works across different assignments.  Rule tests cover the diverse set of
condition types provided by the library, including equality, inequality,
comparison and regular expression operators.  Composite logic such as AND, OR,
NOT and implication is also checked for correct evaluation.

The parsing module is verified with ten dedicated tests. These construct small
machine definitions in the Robot Framework style and feed them to the
\texttt{parse} function. Parsing handles variables, settings, user-defined
keywords and rule declarations. The following snippet from
\texttt{test\_parsing.py} illustrates how conditional actions and rules are
validated:
\begin{lstlisting}[language=python]
machine_text = """*** Machine ***
${VAR1}  any of  a  b
${VAR2}  any of  1  2

${VAR1} == a  ==>  ${VAR2} == 1

Start
  [Actions]
    action1  ==>  End

End
"""
machine = parse(machine_text)
self.assertEqual(len(machine.rules), 1)
self.assertTrue(machine.rules_are_ok(["a", "1"]))
self.assertFalse(machine.rules_are_ok(["a", "2"]))
\end{lstlisting}
Each test confirms that the parser not only instantiates the machine correctly
but also evaluates the resulting rules as expected.

Generation strategies are checked in isolation via fifteen tests. The
\texttt{DepthFirstSearchStrategy} is validated for deterministic path
exploration, action limits and optional target states. The
\texttt{RandomStrategy} is stressed with repeated value generation and random
walks to ensure variability.  Generator unit tests focus on producing Robot
Framework output. They confirm coverage tracking, limit handling and elimination
of duplicate tests.  Together, these unit tests guard each component against
regressions and provide confidence that the smaller building blocks behave as
specified.

Several tests purposefully target boundary scenarios. Examples include empty
machines with no actions, invalid state names and variable identifiers that do
not match the required regular expression. These checks ensure that the parser
and model raise appropriate assertions rather than failing silently. All unit tests rely on
Python's \texttt{unittest} framework. 

\subsection{Integration Tests}

Integration tests combine parsing, strategy execution and output formatting to
validate end-to-end behaviour.  Nine such tests are defined in
\texttt{test\_integration.py}.  They construct complete machine definitions,
parse them into model objects and invoke the generator with different strategies
and configuration options.  The resulting Robot Framework text is inspected for
correct structure, coverage annotations and valid keyword sections.

A representative example is the \texttt{test\_parse\_and\_generate\_simple\_machine}
case.  It builds a small machine containing settings and variables, parses it and
runs the generator while tracking which states and actions were visited.  The
snippet below shows the verification of the produced output:
\begin{lstlisting}[language=python]
all_actions = set()
for state in machine.states:
    all_actions.update(state._actions)

generator.generate(
    machine,
    max_tests=10,
    max_actions=3,
    output=output,
    strategy=DepthFirstSearchStrategy,
    all_actions=all_actions
)
content = output.getvalue()
self.assertIn("*** Settings ***", content)
self.assertIn("TEST COVERAGE INFORMATION", content)
\end{lstlisting}
The test ensures that the pipeline from textual specification to final Robot
Framework file executes without error and that coverage data is embedded in the
output. Coverage counters for states and actions are asserted to appear at the
top of the generated file so that developers can see exactly which parts of the
machine were exercised.

Integration tests are executed using Python's \texttt{unittest}. 

Other integration tests handle more complex scenarios: generating from a machine
with extensive rules, performing random exploration, limiting tests to those that
reach a given target state and running the convenience function
\texttt{transform}. One test even creates temporary files to emulate command-line
usage. These tests confirm that the system works reliably as a whole and can
handle larger inputs without excessive runtime.

Several scenarios also probe error handling. One integration test feeds the
generator a machine whose rules cannot all be satisfied, verifying that test
generation still completes gracefully. Another constructs a large state machine
containing many variables and branches to observe performance and memory
consumption under stress. 

\subsection{Login FSM Example}

The directory \texttt{tests/generator/example/} provides a complete login state
machine and the Robot Framework tests generated by each strategy based on the FSM in figure~\ref{fig:fsm_example}.  The machine is
specified in the following \texttt{.machine} format:
\begin{lstlisting}[language=,caption={Example FSM defined in .machine format}, label={lst:example_machine}]
*** Variables ***
${VALID_USER}         user1
${INVALID_USER}       invalid_user

*** Machine ***
${USER}       any of    ${VALID_USER}      ${INVALID_USER}

Start
  [Actions]
    authenticate user        ==>  Authenticated    when  ${USER} == ${VALID_USER}
    authenticate user        ==>  Failed           when  ${USER} == ${INVALID_USER}

Authenticated
  [Actions]
    logout                   ==>  Start
\end{lstlisting}
When processed by the generator, three strategies yield markedly different test
suites. Table~\ref{tab:strategy_comparison} summarizes the key characteristics
and coverage metrics for each approach.

\begin{table}[ht]
\centering
\caption{Comparison of test generation strategies applied to the example FSM}
\label{tab:strategy_comparison}
\begin{tabular}{|l|c|c|c|c|}
\hline
\textbf{Strategy} & \textbf{Test Cases} & \textbf{States Covered} & \textbf{Actions Covered} & \textbf{Characteristics} \\
\hline
Depth-First & 51 & 3/6 (50\%) & 5/12 (42\%) & 
\begin{tabular}[c]{@{}c@{}}Deterministic sequences\\Systematic exploration\\Misses failure paths\end{tabular} \\
\hline
Random & 51 & 6/6 (100\%) & 12/12 (100\%) & 
\begin{tabular}[c]{@{}c@{}}Stochastic exploration\\Full coverage\\Variable outcomes\end{tabular} \\
\hline
All-Pairs & 5 & 6/6 (100\%) & 12/12 (100\%) & 
\begin{tabular}[c]{@{}c@{}}Variable interactions\\Compact test suite\\Efficient coverage\end{tabular} \\
\hline
\end{tabular}
\end{table}

Table~\ref{tab:detailed_coverage} provides a detailed breakdown of which specific
states and actions are covered by each strategy, highlighting the differences in
exploration patterns.

\begin{table}[ht]
\centering
\caption{Detailed coverage breakdown by generation strategy}
\label{tab:detailed_coverage}
\small
\begin{tabular}{|l|p{3.5cm}|p{3.5cm}|}
\hline
\textbf{Strategy} & \textbf{Covered States} & \textbf{Uncovered States} \\
\hline
Depth-First & Authenticated, Start, Success & End, Error, Failed \\
\hline
Random & All states (6/6) & None \\
\hline
All-Pairs & All states (6/6) & None \\
\hline
\end{tabular}

\vspace{0.5cm}

\begin{tabular}{|l|p{8cm}|p{2cm}|}
\hline
\textbf{Strategy} & \textbf{Key Uncovered Actions} & \textbf{Missing Paths} \\
\hline
Depth-First & 
\begin{tabular}[c]{@{}l@{}}
• perform action (Auth $\to$ Error) \\
• authenticate user (Start $\to$ Failed) \\
• exit system (Failed $\to$ End) \\
• retry authentication (Failed $\to$ Start)
\end{tabular} & Failure scenarios \\
\hline
Random & None (full coverage) & None \\
\hline
All-Pairs & None (full coverage) & None \\
\hline
\end{tabular}
\end{table}

\textbf{Depth-first search} generates long, deterministic sequences that
systematically explore branches until the maximum action limit is reached, but
may miss unreachable states due to its deterministic nature.
\textbf{Random generation} achieves full coverage through stochastic exploration,
visiting all states and actions through variability in path selection.
\textbf{All-pairs} creates the most efficient test suite, achieving complete
coverage with only five tests by focusing on variable interactions and pairwise
combinations.

Overall, the example demonstrates how the generator supports different testing
approaches. Depth-first yields exhaustive paths but may miss unreachable states,
random exploration improves coverage through variability, and all-pairs efficiently
covers variable interactions with few tests.
